\documentclass[11pt]{article}
\newif\ifshowsolutions

%%%%%%%%%%%%%%%%%%%%%%%%%%%%%%%%%%%%%%%%%%%%%%%%%%%%%%%%%%%%%%%%%%%%%%%%%%%%%
% Choose to show solutions (\showsolutionstrue) or not (\showsolutionsfalse).
% (Commenting this out entirely should hide the solutions).
\showsolutionsfalse %true or false
%%%%%%%%%%%%%%%%%%%%%%%%%%%%%%%%%%%%%%%%%%%%%%%%%%%%%%%%%%%%%%%%%%%%%%%%%%%%%

\usepackage[normalem]{ulem}
\usepackage{amsfonts,amsmath,amssymb}
\usepackage{rotating}
\usepackage{color,graphicx}
\usepackage{fullpage,setspace}
\usepackage[colorlinks=false,urlcolor=darkgray,bookmarks=false,pdfborder={0 0 0}]{hyperref}
\usepackage{sgame,multirow,pstricks,egameps}
\usepackage{verbatim}
\usepackage{natbib}
\newcounter{question}
\newtheorem{problem}{Problem}[question]
\newcommand{\newquestion}[1][]{\stepcounter{question}\bigskip{\large \textbf{\underline{Problem \arabic{question}{#1}}}}\bigskip}

\newenvironment{proof}[1][Proof]{\noindent\textbf{#1.} }{\ \rule{0.5em}{0.5em}}
\newenvironment{solution}[1][Solution]{\doublespacing\noindent\textbf{#1 \theproblem}\ \ }{\hfill\rule{0.5em}{0.5em} %don't delete these 2 line breaks:

\bigskip}

\ifshowsolutions
    % do nothing... solutions will be shown.
\else
    % hide the solutions as comments.
    \let\solution=\comment \let\endsolution=\endcomment\bigskip
\fi

\newtheorem{theorem}{Theorem}[question]
\renewcommand{\thetheorem}{\arabic{question}.\Alph{theorem}}
\newtheorem{axiom}[theorem]{Axiom}
\renewcommand{\theaxiom}{\arabic{question}.\Alph{axiom}}
\newtheorem{assumption}[theorem]{Assumption}
\renewcommand{\theassumption}{\arabic{question}.\Alph{assumption}}
\newtheorem{claim}[theorem]{Claim}
\renewcommand{\theclaim}{\arabic{question}.\Alph{claim}}
\newtheorem{conjecture}[theorem]{Conjecture}
\renewcommand{\theconjecture}{\arabic{question}.\Alph{conjecture}}
\newtheorem{definition}[theorem]{Definition}
\renewcommand{\thedefinition}{\arabic{question}.\Alph{definition}}
\newtheorem{example}[theorem]{Example}
\renewcommand{\theexample}{\arabic{question}.\Alph{example}}
\newtheorem{proposition}[theorem]{Proposition}
\renewcommand{\theproposition}{\arabic{question}.\Alph{proposition}}
\newtheorem{remark}[theorem]{Remark}
\renewcommand{\theremark}{\arabic{question}.\Alph{remark}}
\newtheorem{corollary}[theorem]{Corollary}
\renewcommand{\thecorollary}{\arabic{question}.\Alph{corollary}}
\newtheorem{lemma}[theorem]{Lemma}
\renewcommand{\thelemma}{\arabic{question}.\Alph{lemma}}

\newtheorem{case}{Case}[theorem]
\renewcommand{\thelemma}{\arabic{question}.\Alph{theorem}.\alph{case}}

\numberwithin{figure}{question}
\setlength{\parindent}{0mm}
\def\changemargin#1#2{\list{}{\rightmargin#2\leftmargin#1}\item[]}
\let\endchangemargin=\endlist
\definecolor{scarlet}{cmyk}{0,1.00,0.65,0.15}
\definecolor{gray}{cmyk}{0.06,0,0,0.34}
\definecolor{darkgray}{cmyk}{0.07,0,0,0.90}

\begin{document}
\vspace*{-20mm}
\begin{changemargin}{-0.5in}{0.5in}
\footnotesize{
\begin{tabular}{lr}
\multirow{1}{*}{\includegraphics
[height=0.66in,width=0.66in]{graphics/osu_logo.pdf}} & \vspace{-5pt}\textbf{\small{Department of Economics}}\\
                &\color{scarlet}{\rule{6.375in}{0.4pt}}\hspace*{1pt}\\
                & 410 Arps Hall\\
                & 1945 North High Street\\
                & Columbus, OH 43210-1172\\
                & \\
\textbf{ECON 8877} & \textbf{Prof. Paul J. Healy}\\
Fall 2025           & Email: \href{mailto:healy.52@osu.edu}{healy.52@osu.edu}\\
\end{tabular}
}
\end{changemargin}

\begin{center}
\bigskip

{\Large Hypothesis Test Simulation Assignment
\ifshowsolutions
    - Solutions
\fi }

%\textbf{Due by Monday, November 6}$^{\textbf{th}}$\textbf{ 2023 at 4:00 PM}

\smallskip

\end{center}

Please create a PDF showing all your estimated power curves, along with your analyses. Be prepared to present them to the class.

\newquestion[: When is the $t$-test Valid?]

We'll focus on the one-sample, two-sided $t$-test.

Using whatever software you want, you're going to simulate $R$ runs of an experiment (or dataset) with $n$ observations, so $X^r=(X_1^r,\ldots,X_n^r)$ where each $X_i^r\overset{iid}{\sim} F$ for some ``true'' distribution $F_\theta$ (which has parameter vector $\theta$). Our hypotheses are going to be of the form:
$$
    H_0:E_{F_\theta}[X_i]=\mu_0 \hspace{1in} H_1:E_{F_\theta}[X_i]\neq \mu_0,
$$
where $F_\theta$ is the distribution that generates $X_i$. Normally we'd make assumptions about $F_\theta$---like, $F_\theta$ is normal with $\theta=(\mu,\sigma)$---but here we're going to ask ``under what types of $F_\theta$ does the $t$-test have good properties?'' To do so, you'll pick some $F_\theta$ with parameter vector $\theta$, run the simulation, and plot your estimated power curve $\hat{\beta}(\theta)$. Remember that the test is \textit{valid} if $\hat{\beta}(\theta)=0.05$ when $E_{F_\theta}[X_i]=\mu_0$ and \textit{consistent} if $\lim_{n\rightarrow\infty} \hat{\beta}(\theta)=1$ when $E_{F_\theta}[X_i]\neq 0$.

\begin{problem}
    First let's verify the $t$-test works when it should work. Generate the data with $F_\theta$ being Normal with mean $\mu$ and whatever variance you want. Let's start with $n$ somewhat modest, somewhere between 30 and 100; whatever you want. Plot the power curve. Verify the test is valid (at $\mu=\mu_0$).
\end{problem}

\begin{problem}
    Increase $n$ quite a lot and see if it looks like the test is consistent.
\end{problem}

\begin{problem}
    Back to small $n$, increase or decrease the variance. How much is power affected?
\end{problem}

\begin{problem}
    Now let's test the limits of the central limit theorem. Pick some non-Normal distribution for $F_\theta$ and let $\theta_0$ be a (or the) parameter vector where $E_{F_{\theta_0}}[X_i]=\mu_0$. For the power curve, show $\hat{\beta}$ for different values of $E_{F_\theta}[X_i]$ (if the mean isn't an explicit parameter then the $x$-axis of your graph should not be a parameter value, but instead be the mean at that parameter vector). Again, start with modest $n$ and check if the test is valid.
\end{problem}

\begin{problem}
    Slowly ramp up $n$. Does it look to be asymptotically valid? Asymptotically consistent? Is $n\geq 30$ the magic threshold?
\end{problem}

\begin{problem}
    Now try to construct a really bimodal distribution and repeat the above exercise. What kind of $n$ do we need to trust the $t$-test?
\end{problem}

Now try the two-sample $t$-test. So you'll generate two samples, $Y^1$ with $n$ observations and $Y^2$ with $m$ observations. The sample $Y^1$ is generated from $F$ and $Y^2$ is generated from $G$. The hypotheses are:
$$
    H_0:E_F[Y_i^1]=E_G[Y_i^2] \hspace{1in} H_1:E_F[Y_i^1]\neq E_G[Y_i^2].
$$
When plotting power curves, the $x$-axis should be the (true) difference in means.

\begin{problem}
    Start with $F$ and $G$ both being Normal with the same variance, and differing only in their mean. Let $\theta=\mu_G-\mu_F$. Verify the test works as advertised.
\end{problem}

\begin{problem}
    Now let $F$ and $G$ be Normal but with different variances. Use the Welch's version of the $t$-test that assumes unequal variances. Does it work as advertised?
\end{problem}

\begin{problem}
    Suppose again $F$ and $G$ are Normal with the same variance, but you run Welch's version of the test. Is it still valid? Do you lose a lot of power by using this test (as opposed to the version that assumes equal variances)?
\end{problem}

\newquestion[: Which Test is Best?]

Consider the same two-sample (with potentially different sample sizes), two-sided testing situation as the the previous question. Now we're going to compare:
\begin{enumerate}
    \item Original (pooled-variance) $t$-test
    \item Welch's $t$-test
    \item Wilcoxon ranksum test
    \item Permutation test for difference in means
    \item Variance-corrected permutation test for difference in means \citep{ChungRomano2013}
\end{enumerate}

\begin{problem}
    Start with the nicest world: both distributions are Normal with the same variance. Plot power curves for all five tests. Are they all valid? Which gives the greatest power?
\end{problem}

\begin{problem}
    Now try Normal distributions but with unequal variances. Are all tests still valid? Which gives the greatest power?
\end{problem}

\begin{problem}
    Now try some very non-Normal distributions (\textit{e.g.}, bimodal distributions). Which test wins now?
\end{problem}

\newquestion[: Using Real Data]

\begin{problem}
    Find a real dataset with two samples where we're interested in differences in means. Do NOT plot histograms of the data yet. Just run the five tests above and report the results.
\end{problem}

\begin{problem}
    Now plot histograms of the two samples. Do they look Normal? Based on what you learned from your simulations, can you predict which test is actually best to use here?
\end{problem}

\bibliographystyle{elsarticle-harv}
\bibliography{pjbib}

\end{document}
